We test whether a transition-dependent finite-group cocycle can repair the
arithmetic overproduction of the tensor-subset full shift.  Relabeling
naturality and restriction compatibility classify every source-locked local rule
by the three incidence counts
$(|S\setminus T|,|S\cap T|,|T\setminus S|)$; this leaves
$\binom{n+3}{3}-(2n+1)$ types on $n$ atoms.  We then assign
$r=(12)$ to strict refinements, $t=(23)$ to strict coarsenings, and the
identity otherwise.  The resulting $S_3$ extension is a genuine deck
symmetry of the same two-block symbolic object and is not cohomologous to a
one-letter cardinality clock.  Its two-atom trivial and sign determinants
are both $(1-x)(1-y)$, whereas its standard block is exactly
\[
(1-x)^2(1-y)^2+3xy(x+y)(xy+1)(x+y-1).
\]
The relative trace logarithm leaks at $x^2y,xy^2,x^2y^2$ with coefficients
$-3,-3,-6$.  A primitive four-edge word has commutator holonomy $(rt)^2$
and an edge-separated standard-character gap of three, so the nonabelian
signal cannot be dismissed as a coboundary or an unmarked-orbit
cancellation.  Exhaustive two-atom searches in $S_3,D_4,Q_8$ found that
every all-irrep-clean table belongs to the counting-gauge class; this is
finite evidence, not a general theorem.  The infinite nontrivial block is
trace class for $\RePart s>2$, while the trivial block is already defined
for $\RePart s>1$.  Thus transition holonomy is dynamically genuine but
inventory-blind: it preserves the scalar Euler factor yet cannot select
prime or prime-power primitives.  The strict outcome is
\texttt{ROUTE\_A\_REJECTED}; no claim about Riemann zeros is made.
